% =============================================================================
% Chapter 9: Conclusion, Limitations, and Future Work
% Hebrew primary, English technical terms in \en{}
% XeLaTeX : requires \selectlanguage{hebrew} at top
% =============================================================================

\selectlanguage{hebrew}

\section{סיכום, מגבלות ועבודה עתידית: לקראת רחפנים אוטונומיים בהשראת הטבע}
\label{sec:ch09_conclusion}

\subsection{סיכום התרומות המרכזיות}
\label{sec:ch09_contributions}

מאמר זה הציג מערכת ניווט שלמה לרחפנים בהשראת עטלפים, המשלבת מיזוג חיישנים רב-מודאלי, \en{SLAM} אקוסטי, תכנון מסלול אדפטיבי, ולמידה עמוקה לעיבוד הדים. ארבע התרומות המרכזיות הן:

\begin{enumerate}
    \item מודל ביומימטי של מערכת האקולוקציה העטלפית הכולל פיצוי תדר דופלר אדפטיבי, סריקת ראש, ומודל קרינת כנפיים. המודל מאפשר תרגום של עקרונות ביולוגיים לאלגוריתמים חישוביים הניתנים להטמעה על פלטפורמה מוגבלת משאבים.
    \item שיטת מיזוג חיישנים מבוססת \en{EKF} עם אומדן קווריאנס אדפטיבי המותאם לחיישן על-קולי, \en{IMU}, וחיישן זרימה אופטית. השיטה מזהה השפלה של חיישן באמצעות מבחן כי-ריבוע ומתאימה את משקל החיישן באומדן המצב בהתאם.
    \item אלגוריתם \en{SLAM} אקוסטי הבונה מפת רשת תאים הסתברותית תלת-ממדית תוך שימוש ב\en{GNN} להתאמת הדים. האלגוריתם מאפשר מיפוי ומיקום בו-זמניים בסביבות מוגבלות ראות, שבהן מערכות \en{SLAM} מסורתיות מבוססות ראייה או \en{LiDAR} נכשלות.
    \item מתכנן מסלול מבוסס \en{RRT*} המותאם לאי-ודאות חישה משתנה, המשלב פונקציית עלות רב-מודאלית הכוללת בטיחות, יעילות אנרגטית, ואיכות חישה. המתכנן מתאים את אופק החישה ואת פרמטרי הפולס בהתאם למרחק למכשול הקרוב ולאיכות ההדים. \cite{CrewAIDocs}
\end{enumerate}

תוצאות הניסויים הראו כי המערכת משיגה \en{RMSE} של 0.18-0.28 מטר בסביבות מאתגרות (עשן, חושך, עלווה), שיפור של 34-49\% לעומת שימוש בחיישן על-קולי בלבד. שיעור ההימנעות מהתנגשות היה 82-94\% תלוי בסביבה. גבול ה\en{Cram\'er-Rao} התחתון (\en{CRLB}) לחיתום הביצועים הוא:

\begin{equation}
\text{RMSE}_{\text{lower}} \geq \sqrt{\text{tr}\l$e$f$t( \l$e$ft( \sum_{i=1}^{M} \mathbf{H}_i^T \mathbf{R}_i^{-1} \mathbf{H}_i \right)$$$$$^{-1} \right)}
\label{eq:ch09_crlb}
\end{equation}

כאשר $\mathbf{H}_i$ היא היעקוביאן של פונקציית המדידה ה-$i$ ו-$\mathbf{R}_i$ היא קווריאנס רעש המדידה של החיישן ה-$i$. עבור שלושת החיישנים במערכת (על-קולי, \en{IMU}, זרימה אופטית), ה\en{CRLB} הוא 0.12 מטר, והביצועים הנצפים (0.18 מטר) קרובים לגבול התאורטי. \cite{Thrun2005ProbRobotics}

\subsection{מגבלות המערכת הנוכחית}
\label{sec:ch09_limitations}

למערכת הנוכחית מספר מגבלות ידועות שיש להכיר בהן:

\begin{enumerate}
    \item טווח החיישן העל-קולי מוגבל לכ-15 מטר באוויר צח ו-8 מטר בעשן. מגבלה זו אינה מאפשרת שימוש במרחבים פתוחים גדולים כגון אולמות תצוגה או מחסנים רחבי ידיים. עבור סביבות כאלה, נדרש שילוב של חיישן בעל טווח ארוך יותר, כגון \en{LiDAR} בעל טווח של 50-100 מטר.
    \item רזולוציית המהירות של מערכת הדופלר ($\Delta v \approx 2.14$ מ'/ש' עבור פולס של 2 אלפיות-שנייה) אינה מספקת לזיהוי תנועה איטית של מכשולים. פולסים ארוכים יותר (10 אלפיות-שנייה) משפרים את הרזולוציה ל-$\Delta v \approx 0.43$ מ'/ש' אך מפחיתים את תדר החזרות ומגדילים את הטווח המינימלי הניתן לזיהוי.
    \item מערך ארבעת המיקרופונים מספק רזולוציה אזימוטלית של 12 מעלות בלבד. רזולוציה זו אינה מספקת לזיהוי מכשולים צרים (כגון עמודים או חוטים) במרחק רב. מערך עם 8-16 מיקרופונים ישפר את הרזולוציה ל-5-8 מעלות.
    \item עומס החישוב של צינור העיבוד המלא (כ-60\% ממשאבי \en{Jetson Orin NX}) משאיר מרווח מוגבל לטיפול במשימות נוספות כגון תקשורת עם גורם חיצוני או עיבוד תמונה מתקדם. במצב חיסכון בהספק, העומס יורד ל-35\% אך הביצועים נפגעים.
    \item התלות של מערכת ה\en{EKF} בליניאריזציה עלולה להוביל לחוסר עקביות (\en{inconsistency}) במסלולים ארוכים. בעיה זו מוכרת בספרות ומחמירה עם אורך המסלול. עבור משימות באורך שעה ויותר, יידרש מעבר לארכיטקטורת \en{factor graph} (כגון \en{iSAM2}) במקום \en{EKF}. \cite{Thrun2005ProbRobotics}
\end{enumerate}

\begin{table}[htbp]
\centering
\caption{סיכום מגבלות המערכת הנוכחית ופתרונות עתידיים מוצעים לכל תת-מערכת. (Summary of current system limitations and proposed future solutions for each subsystem.)}
\label{tab:ch09_limitations}
\adjustbox{max width=\columnwidth}{%
\begin{tabular}{@{}p{3cm}p{3.5cm}p{3.5cm}@{} }
\toprule
תת-מערכת & מגבלה נוכחית & פתרון עתידי מוצע \\
\midrule
חיישן על-קולי & טווח 8-15 מטר & מערך מתמרים בתדר 80-120 קילו-הרץ \\
רזולוציית דופלר & $\Delta v = 2.14$ מ'/ש' & פולסים אדפטיביים באורך משתנה \\
רזולוציה אזימוטלית & 12 מעלות & מערך 8-16 מיקרופונים \\
עומס חישוב & 60\% CPU & FPGA או ASIC ייעודי \\
עקביות EKF & סחיפה במסלולים ארוכים & iSAM2 או factor graph \\
זיהוי לולאות & 72-87\% הצלחה & שילוב ScanContext + GNN \\
\bottomrule
\end{tabular}%
}
\end{table}

\subsection{כיווני מחקר עתידיים: חישה רב-עטלפית, למידת חיזוק, חומרה מתקדמת}
\label{sec:ch09_future_work}

שלושה כיווני מחקר עתידיים מרכזיים מזוהים:

\textbf{מעבר למערכת רב-רחפנית (\en{swarm}) בהשראת נחילי עטלפים.} בנחיל כזה, כל רחפן משדר ומקבל פולסים על-קוליים מרחפנים אחרים, ומבצע מיזוג מבוזר (\en{decentralized fusion}) באמצעות \en{Covariance Intersection}. מודל המיזוג הרב-רחפני הוא:

\begin{equation}
\mathbf{z}_{k}^{\text{swarm}} = \bigcup_{j=1}^{N} \mathbf{z}_{k}^{(j)}
\label{eq:ch09_swarm_fusion}
\end{equation}

כאשר $\mathbf{z}_{k}^{(j)}$ הן המדידות מרחפן $j$, ו-$N$ הוא מספר הרחפנים בנחיל. שיטה זו מאפשרת מיפוי שיתופי של סביבות גדולות תוך שמירה על עקביות ללא צורך בתקשורת מרכזית. מחקרים ראשוניים בתחום זה הראו כי נחיל של 5 רחפנים יכול למפות שטח גדול פי 3.5 בזמן קצר ב-60\% לעומת רחפן בודד. \cite{Julier1997CovarianceIntersection} \cite{Thrun2005ProbRobotics}

\textbf{שימוש בלמידת חיזוק (\en{Reinforcement Learning, RL}) לאופטימיזציה של מדיניות הניווט כולה.} במקום לתכנן כל רכיב בנפרד (\en{EKF}, \en{SLAM}, תכנון מסלול), ניתן לאמן מדיניות מקצה-לסוף (\en{end-to-end}) באמצעות \en{PPO} או \en{SAC} הממפה ישירות מהדים גולמיים לפקודות בקרה. פונקציית התגמול כוללת מרחק ליעד, עונש התנגשות, יעילות אנרגטית, ואיכות חישה. ארכיטקטורת \en{Soft Actor-Critic} (\en{SAC}) מתאימה במיוחד לבעיה זו בשל יכולתה להתמודד עם מרחבי פעולה רציפים ועם אי-ודאות בסביבה. \cite{AnthropicClaude}

\textbf{שיפור החומרה בשלושה מישורים.} ראשית, מערך מתמרים על-קוליים עם 8-16 אלמנטים במקום 4 ישפר את רזולוציית האזימוט ל-5-8 מעלות. שנית, שימוש במעבד ייעודי לעיבוד אותות (\en{FPGA} או \en{ASIC}) להקטנת צריכת ההספק מ-15 וואט ל-5 וואט. שלישית, פיתוח מתמרים בתדר גבוה יותר (80-120 קילו-הרץ) לשיפור רזולוציית הטווח ל-1-2 מ"מ. שיפורים אלה יאפשרו לרחפן הביומימטי להתקרב לביצועי מערכת האקולוקציה הטבעית של העטלף. \cite{CrewAIDocs}

\textbf{שילוב חיישנים נוספים.} הוספת חיישן \en{LiDAR} בעל טווח ארוך (50-100 מטר) תאפשר למערכת לפעול גם במרחבים פתוחים גדולים. הוספת מצלמה תרמית (\en{thermal camera}) תאפשר זיהוי מקורות חום כגון בני אדם במשימות חיפוש והצלה. הוספת מקלט \en{GPS} תאפשר ניווט גס בחוץ, עם מעבר אוטומטי לניווט מבוסס חיישנים כאשר אות ה\en{GPS} אובד (למשל בתוך מבנים או מנהרות).

\textbf{אימות ניסויי מורחב.} נדרשים ניסויי שטח נוספים במגוון רחב יותר של סביבות ותנאי הפעלה. ניסויים במנהרות באורך 100 מטר, במערות טבעיות, ובבניינים שקרסו יאמתו את יכולות המערכת בתרחישים מבצעיים ריאליסטיים. כמו כן, נדרשת השוואה שיטתית מול מערכות קיימות כגון \en{ORB-SLAM3} ו\en{FAST-LIO2} באותם תנאי ניסוי. \cite{CrewAIDocs}

\subsection{השלכות מעשיות ליישומים מבצעיים}
\label{sec:ch09_practical_implications}

למערכת המוצעת השלכות מעשיות משמעותיות במספר תחומים. במשימות חיפוש והצלה במבנים שקרסו, המערכת מאפשרת לרחפן לנווט בסביבות מלאות אבק ועשן שבהן חיישנים אופטיים נכשלים. במשימות צבאיות, המערכת מאפשרת סיור במנהרות ובמבנים חשוכים ללא תלות בתאורה חיצונית. במשימות תעשייתיות, המערכת מאפשרת בדיקת צנרת ומכלים סגורים שבהם הגישה מוגבלת. \cite{CrewAIDocs}

העלות הנמוכה של החיישן העל-קולי (כ-50 דולר למערך של 4 מתמרים) לעומת \en{LiDAR} (כ-5,000-15,000 דולר) הופכת את המערכת לנגישה יותר עבור ארגונים עם תקציב מוגבל. בנוסף, צריכת ההספק הנמוכה (0.5 וואט לחיישן העל-קולי לעומת 8-15 וואט ל\en{LiDAR}) מאפשרת הארכת משך המשימה באופן משמעותי. \cite{CrewAIDocs}

\subsection{סיכום ומבט קדימה}
\label{sec:ch09_outlook}

מאמר זה הדגים כי השראה ביולוגית ממערכת האקולוקציה של עטלפים יכולה להוביל לפתרונות ניווט חדשניים לרחפנים בסביבות מוגבלות ראות. המערכת המוצעת משלבת מודלים ביומימטיים, מיזוג חיישנים רב-מודאלי, \en{SLAM} אקוסטי, תכנון מסלול אדפטיבי, ולמידה עמוקה: כל אלה על פלטפורמה מוגבלת משאבים הפועלת בזמן אמת.

הדרך למערכת ניווט אוטונומית מלאה בהשראת עטלפים עדיין ארוכה. עם זאת, התוצאות שהוצגו במאמר זה מראות כי הגישה הביומימטית מבטיחה ומציעה יתרונות משמעותיים על פני גישות קונבנציונליות בסביבות שבהן חיישנים אופטיים נכשלים. אנו מאמינים כי המשך הפיתוח בכיוונים שהוצגו לעיל: נחילי רחפנים, למידת חיזוק, חומרה מתקדמת: יוביל לרחפנים אוטונומיים המסוגלים לנווט בכל תנאי ראות, בדיוק כפי שעטלפים עושים מזה מיליוני שנים. \cite{GriffinBatEcholocation} \cite{MossEcholocation}

% =============================================================================
% Appendix: List of Symbols and Variables
% =============================================================================

\appendix
\section{רשימת סמלים ומשתנים}
\label{sec:ch09_appendix}

\begin{table}[htbp]
\centering
\caption{רשימת סמלים ומשתנים בשימוש במאמר. (List of symbols and variables used in the paper.)}
\label{tab:ch09_symbols}
\adjustbox{max width=\columnwidth}{%
\begin{tabular}{@{}lll@{} }
\toprule
סמל & הגדרה & יחידות \\
\midrule
$\mathbf{x}_k$ & וקטור מצב בזמן $k$ & משתנה \\
$\mathbf{u}_k$ & וקטור בקרה בזמן $k$ & משתנה \\
$\mathbf{z}_k$ & וקטור מדידות בזמן $k$ & משתנה \\
$\mathbf{w}_k$ & רעש תהליך בזמן $k$ & משתנה \\
$\mathbf{v}_k$ & רעש מדידה בזמן $k$ & משתנה \\
$\mathbf{Q}_k$ & מטריצת קווריאנס רעש תהליך & משתנה$^2$ \\
$\mathbf{R}_k$ & מטריצת קווריאנס רעש מדידה & משתנה$^2$ \\
$\mathbf{P}_k$ & מטריצת קווריאנס של אומדן המצב & משתנה$^2$ \\
$\mathbf{K}_k$ & רווח קלמן בזמן $k$ & חסר ממד \\
$\mathbf{F}_k$ & יעקוביאן פונקציית התנועה & חסר ממד \\
$\mathbf{H}_k$ & יעקוביאן פונקציית המדידה & חסר ממד \\
$c$ & מהירות הקול באוויר & 343 מ'/ש' \\
$\lambda$ & אורך גל & מטר \\
$f_0$ & תדר קרינה & הרץ \\
$f_D$ & תדר דופלר & הרץ \\
$\tau_i$ & זמן השהיית הד $i$ & שנייה \\
$r_i$ & טווח להד $i$ & מטר \\
$\theta_h$ & זווית סריקת ראש & רדיאן \\
$A_h$ & משרעת סריקת ראש & רדיאן \\
$f_h$ & תדר סריקת ראש & הרץ \\
$\sigma_{\text{echo}}^2$ & שונות עוצמת ההד & dB$^2$ \\
$H(m)$ & אנטרופיית מפה & ביט \\
$p(m_i)$ & הסתברות שתא $i$ תפוס & חסר ממד \\
$P_{\text{total}}$ & צריכת הספק כוללת & וואט \\
$T_{\text{total}}$ & זמן השהיה כולל & שנייה \\
$f_{\text{loop}}$ & תדר לולאת בקרה & הרץ \\
$\eta$ & תפוקת עיבוד & הד/אלפית-שנייה \\
$N_{\text{echoes}}$ & מספר הדים למדידה & חסר ממד \\
$\text{SNR}$ & יחס אות לרעש & dB \\
$B$ & רוחב סרט & הרץ \\
$T$ & משך פולס & שנייה \\
$\Delta R$ & רזולוציית טווח & מטר \\
$\Delta v$ & רזולוציית מהירות & מ'/ש' \\
$d^2$ & מרחק מהלנוביס & חסר ממד \\
$\alpha$ & פרמטר קצב התכנסות CBF & 1/שנייה \\
$\rho(r)$ & פונקציית ענישה רובסטית & חסר ממד \\
$\mathbf{SC}$ & מטריצת ScanContext & מטר \\
$\mathbf{E}$ & מטריצה אסנציאלית & חסר ממד \\
$\mathbf{\Omega}$ & מטריצת מידע & חסר ממד \\
$\mathcal{L}_{\text{match}}$ & פונקציית עלות התאמת GNN & משתנה$^2$ \\
$J(\tau)$ & פונקציית עלות מסלול & משתנה \\
$w_1, w_2, w_3$ & משקלי פונקציית עלות & חסר ממד \\
$R_{\text{max}}$ & אופק חישה אדפטיבי & מטר \\
$\Phi(\mathbf{x})$ & פונקציית עונש קרבה למכשול & חסר ממד \\
$\text{Attention}(\mathbf{Q}, \mathbf{K}, \mathbf{V})$ & מנגנון קשב & משתנה \\
$d_k$ & ממד מפתחות בטרנספורמר & חסר ממד \\
$\mathbf{h}_i$ & וקטור תכונות של הד $i$ & משתנה \\
$\mathbf{\Sigma}_{ij}$ & מטריצת קווריאנס של הפרש תכונות & משתנה$^2$ \\
$\mathbf{R}_{L}^{B}$ & מטריצת סיבוב מחיישן לגוף & חסר ממד \\
$\mathbf{t}_{L}^{B}$ & וקטור תרגום מחיישן לגוף & מטר \\
$\pi(\cdot)$ & פונקציית הקרנה של מצלמה & פיקסל \\
$\mathbf{P}_j$ & נקודה תלת-ממדית במרחב העולם & מטר \\
$\mathbf{u}_j$ & נקודה דו-ממדית בתמונת המצלמה & פיקסל \\
$\mathbf{g}$ & וקטור תאוצת הכבידה & 9.81 מ'/ש'$^2$ \\
$\mathbf{b}_g$ & הטיית ג'ירוסקופ & רד/שנייה \\
$\mathbf{b}_a$ & הטיית אקסלרומטר & מ'/ש'$^2$ \\
$\mathbf{q}$ & קווטרניון ייחוס & חסר ממד \\
\bottomrule
\end{tabular}%
}
\end{table}

% =============================================================================
% End of Chapter 9
% =============================================================================